\section{Exact all-phase cone obstruction}
\label{sec:cone}

The feasibility problems \eqref{eq:length-feasibility} and
\eqref{eq:transfer-feasibility} are homogeneous systems of strict rational
inequalities.  We avoid numerical optimization entirely.  Every failure is
recorded by one of two short integer certificates.

\begin{definition}[Certificate types]
\label{def:certificate-types}
Let $\mathcal R$ be the finite collection of all coordinate rows of all
required prefix matrices in a given phase.
\begin{enumerate}[label=(\roman*)]
\item A \emph{nonpositive-row certificate} is a row
  $r\in\mathcal R\setminus\{0\}$ with $r_j\leq0$ for every $j$.
\item A \emph{positive-row-combination certificate} consists of rows
  $r_1,\dots,r_m\in\mathcal R$ and integers $c_1,\dots,c_m>0$ satisfying
  \[
    \sum_{j=1}^m c_jr_j=0.
  \]
\end{enumerate}
\end{definition}

\begin{lemma}[Exact infeasibility]
\label{lem:exact-infeasibility}
Either certificate in \Cref{def:certificate-types} proves that no
$x\in\mathbb R_{>0}^4$ makes every required row strictly positive.
\end{lemma}

\begin{proof}
For type (i), $r\cdot x\leq0$ for every $x>0$, contradicting the required
strict inequality.  For type (ii), if each $r_j\cdot x>0$, then
$\sum_jc_j(r_j\cdot x)>0$, whereas the displayed row identity makes the same
quantity zero.
\end{proof}

This is the elementary strict-cone form of a Farkas alternative.  Because the
matrices, rows, and coefficients are integral, the conclusion has no
tolerance or rounding parameter.

\subsection{The first length-six word}

\begin{theorem}[All phases of $C_1$ fail]
\label{thm:C1-all-phases}
For the word $C_1$ in \eqref{eq:C1}, all six cyclic phases are infeasible for
the forward-length action \eqref{eq:length-prefix}, and all six cyclic phases
are infeasible for the transfer action \eqref{eq:transfer-prefix}.
\end{theorem}

\begin{proof}
In forward-length phase zero, the fourth coordinate of the first prefix and
the second coordinate of the fourth prefix have rows
\begin{equation}
 r=(0,-1,0,1),\qquad -r=(0,1,0,-1).
 \label{eq:C1-opposite-rows}
\end{equation}
Both dot products would have to be strictly positive on a feasible $x$, but
their sum is zero.  This is a type-(ii) certificate with coefficients $1,1$.
For each of the other five rotations, exact chronological replay produces a
certificate of one of the two types in \Cref{def:certificate-types}.
Reversing each raw rotation and replaying \eqref{eq:transfer-prefix} gives six
further integer certificates.  The full descriptors, including step and
coordinate, are generated by the protocol in \Cref{app:certificates}; each
can be checked by integer matrix multiplication and a row equality.  Hence
all twelve strict systems are infeasible.
\end{proof}

The explicit pair \eqref{eq:C1-opposite-rows} is not a proxy for the other
phases.  Cyclic rotation changes the collection of prefix inequalities, so
the all-phase conclusion relies on six separately reconstructed certificates
for each action.

\subsection{The second length-six word}

\begin{theorem}[All phases of $C_2$ fail]
\label{thm:C2-all-phases}
For the word $C_2$ in \eqref{eq:C2}, all six cyclic phases are infeasible for
the forward-length action, and all six cyclic phases are infeasible for the
transfer action.
\end{theorem}

\begin{proof}
In forward-length phase zero, the third coordinate row at step three is
\begin{equation}
 (-1,0,0,0).
 \label{eq:C2-negative-row}
\end{equation}
It is strictly negative on every $x\in\mathbb R_{>0}^4$, so this phase has a
type-(i) certificate.  The remaining five length rotations and the six
contravariant transfer rotations each have an independently reconstructed
integer certificate.  Applying \Cref{lem:exact-infeasibility} to the twelve
descriptors proves the claim.
\end{proof}

\subsection{The length-twenty-four relation}

\begin{theorem}[All phases of $W_{24}$ fail]
\label{thm:W24-all-phases}
Every one of the twenty-four cyclic phases of $W_{24}$ is infeasible for the
forward-length action, and every one is infeasible for the transfer action.
For each action, the canonical certificate census consists of fifteen
nonpositive-row certificates and nine positive-dependence certificates.
\end{theorem}

\begin{proof}
For the forward-length action, phase zero has at step eight the required
second-coordinate row
\begin{equation}
 (-11430,-460520,-3353,-456200),
 \label{eq:W24-length-row}
\end{equation}
which is strictly negative on the positive orthant.  Exact phase enumeration
then gives fifteen type-(i) and nine type-(ii) certificates.

For the transfer action, phase zero begins, in application order, with
$A^{-1},C,B^{-1}$.  At the third prefix, its second-coordinate row is
\begin{equation}
 (-984333,-498163,-999116,-479060).
 \label{eq:W24-transfer-row}
\end{equation}
This is again a type-(i) certificate.  The other twenty-three phases have
exact descriptors reconstructed after reversing their raw path rotations.
Their canonical census is likewise fifteen type-(i) and nine type-(ii)
certificates.
All computations use integer inverses already certified to be unimodular; no
floating-point feasibility solver enters.  Applying
\Cref{lem:exact-infeasibility} phase by phase proves the result.
\end{proof}

Every phase still has endpoint product $I_4$.  Thus
\Cref{thm:W24-all-phases} does not destroy the relation; it shows that no
choice of cyclic basepoint turns that relation into a positive-domain orbit
under either of the two tested dynamical conventions.

\subsection{A control that distinguishes covariant homology}

\begin{proposition}[Raw homology controls]
\label{prop:raw-homology}
For the raw covariant recurrence \eqref{eq:homology-prefix}, the initial
vectors
\begin{equation}
 x_{C_1}=(1,2,1,1)^{\mathsf T},\qquad
 x_{C_2}=(1,1,3,1)^{\mathsf T}
 \label{eq:raw-witnesses}
\end{equation}
have strictly positive coordinates at every prefix of $C_1$ and $C_2$,
respectively.  Shifting the corresponding closed integer trajectories
supplies a positive witness in every cyclic phase.  The analogous raw
covariant feasibility problem for $W_{24}$ is infeasible.
\end{proposition}

\begin{proof}
Multiply the exact chronological matrices successively according to
\eqref{eq:homology-prefix} and evaluate them on the two integer vectors in
\eqref{eq:raw-witnesses}.  Every displayed coordinate is a positive integer,
and the final vector equals the initial one.  A cyclic phase begins at the
corresponding point of this closed trajectory, proving phasewise positivity.
The $W_{24}$ failure has an exact row certificate of the same form as
\Cref{def:certificate-types}.
\end{proof}

\begin{remark}[What the control proves]
The first statement of \Cref{prop:raw-homology} does not rescue either word as
an AGY orbit.  Its role is adversarial: the same word can pass a covariant
positivity calculation while failing both the genuine length and transfer
calculations.  This detects precisely the transpose/order mistake that would
otherwise turn a homology zigzag into a false length orbit.
\end{remark}

\begin{corollary}[Formal-cycle boundary]
\label{cor:formal-boundary}
The three words are exact identity-holonomy cycles of the finite symmetric
path model, but they are not positive-domain periodic orbits for either
frozen Rauzy/AGY action tested above.  Identity holonomy alone is therefore
insufficient to promote their finite trace-log contribution to geometric AGY
periodic data.
\end{corollary}
